\chapter{Fundamentação Teórica}
\label{cap:fundamentacao-teorica}

Este capítulo apresenta os conceitos e dados que embasam o desenvolvimento do trabalho. A fundamentação teórica está organizada em seis seções principais: primeiro, contextualiza-se a pecuária no Brasil, com foco na produção familiar; em seguida, discute-se sua importância econômica; depois, abordam-se os principais desafios enfrentados pelos pequenos produtores; logo após, são exploradas as tecnologias já presentes no campo; na sequência, analisa-se o avanço tecnológico no agronegócio como um todo; por fim, são apresentados os sistemas existentes voltados à gestão da pecuária. O objetivo é demonstrar a relevância do tema e justificar a necessidade de soluções tecnológicas acessíveis para a realidade da agricultura familiar.

\section{Pecuária no Brasil}
\label{sec:fundamentacao-teorica-pecuaria-no-brasil}

A pecuária é uma atividade essencial no contexto agropecuário brasileiro, destacando-se pelo expressivo número de estabelecimentos rurais dedicados a essa prática. O Brasil possui cerca de 5 milhões de estabelecimentos rurais, dos quais 73\% são classificados como agricultura familiar, englobando pequenas e médias propriedades que criam diversas espécies animais, como bovinos, suínos, caprinos, ovinos, equinos e aves \cite{centroAgro2017}.

Em 2023, o efetivo bovino atingiu 238,6 milhões de cabeças, o maior da série histórica iniciada em 1974, evidenciando a relevância e a expansão do setor pecuário no país \cite{agenciaBrasil2023}. Apesar do crescimento do rebanho, a produção de leite apresentou uma redução de 1,6\% no ano anterior, totalizando 34,6 bilhões de litros, fato atribuído pelo IBGE ao aumento dos custos de criação para pequenos produtores.

Além da bovinocultura, outras atividades pecuárias também apresentaram crescimento significativo. O Brasil encerrou 2022 com 12,4 milhões de caprinos, um aumento de 3,9\%, e 21,5 milhões de ovinos, crescimento de 4,7\%, sendo o Nordeste a região líder na criação dessas espécies \cite{agenciaBrasil2023}. A produção de suínos atingiu 44,4 milhões de animais, representando um aumento de 4,3\% em relação a 2021, com destaque para a Região Sul, que concentrou 51,9\% do total nacional, sendo Santa Catarina o estado líder, com 22,1\% do rebanho nacional \cite{agenciaBrasil2023}.

A avicultura também apresentou resultados expressivos. A produção de galináceos, que inclui galinhas e aves para corte, alcançou 1,6 bilhão de cabeças, representando um aumento de 3,8\%. A produção de ovos de galinha foi recorde em 2022, com 4,9 bilhões de dúzias, 1,3\% a mais que no ano anterior, gerando um valor de produção de R\$ 26,1 bilhões, alta de 19,1\% em relação a 2021 \cite{agenciaBrasil2023}.

Na apicultura, a produção de mel em 2022 chegou a 61 mil toneladas, recorde da pesquisa do IBGE, com crescimento de 9,5\% em relação a 2021. O Nordeste foi responsável por 38,7\% da produção nacional, com destaque para os estados do Rio Grande do Sul e Paraná, que lideraram o ranking estadual com 14,8\% e 14,2\% do total nacional, respectivamente \cite{agenciaBrasil2023}.

A aquicultura também se destacou, com a produção de peixes atingindo 617,3 mil toneladas em 2022, um aumento de 6\% na atividade e de 16,4\% no valor de produção, que chegou a R\$ 5,7 bilhões. A Região Sul concentrou 35,8\% do total nacional, com o Paraná sendo o principal estado produtor, responsável por 27,1\% da piscicultura brasileira \cite{agenciaBrasil2023}.

Esses dados evidenciam a diversidade e a importância das diferentes atividades pecuárias na economia brasileira, especialmente no contexto da agricultura familiar, que desempenha papel crucial na produção de alimentos e na geração de renda no meio rural.

\section{Importância Econômica da Pecuária}
\label{sec:fundamentacao-teorica-importancia-economica}

A pecuária brasileira é um dos pilares da economia nacional, contribuindo significativamente para o \gls{PIB} do agronegócio e para a geração de empregos. A cadeia da bovinocultura de corte, por exemplo, representa cerca de 30\% do PIB do agronegócio, sendo o maior segmento exportador do país e responsável por aproximadamente 37\% dos empregos diretos e indiretos no setor. Em 2023, o valor da produção pecuária e da aquicultura alcançou R\$ 122,4 bilhões, com crescimento de 5,4\% em relação ao ano anterior, evidenciando a força econômica do setor \cite{importanciaBovinocultura}. Além disso, a pecuária familiar desempenha papel crucial na segurança alimentar, fornecendo produtos essenciais para o consumo interno.

É importante destacar que, além da bovinocultura, outras espécies animais desempenham papéis significativos na economia brasileira. Em 2022, o rebanho de suínos atingiu 42,3 milhões de cabeças, representando um aumento de 3,2\% em relação ao ano anterior. A produção de carne suína é especialmente relevante nas regiões Sul e Sudeste, com destaque para os estados de Santa Catarina, Paraná e Minas Gerais. A avicultura também apresentou crescimento, com o efetivo de frangos de corte alcançando 1,5 bilhão de cabeças, um aumento de 2,8\% em comparação a 2021. A produção de ovos de galinha totalizou 4,8 bilhões de dúzias, consolidando o Brasil como um dos principais produtores mundiais. No que tange à caprinocultura e à ovinocultura, o país encerrou 2022 com 12,4 milhões de caprinos, um aumento de 3,9\%, e 21,5 milhões de ovinos, crescimento de 4,7\%, sendo o Nordeste a região líder na criação dessas espécies \cite{agenciaBrasil2023}.

Esses dados evidenciam a diversidade e a importância das diferentes cadeias produtivas dentro da pecuária brasileira, ressaltando a necessidade de políticas públicas e investimentos que considerem a pluralidade do setor para promover seu desenvolvimento sustentável e inclusivo.

\section{Desafios Enfrentados na Gestão da Pecuária Familiar}
\label{sec:fundamentacao-teorica-desafios-pecuaria-familiar}

Apesar da importância econômica, as pequenas e médias propriedades enfrentam desafios significativos na gestão das atividades pecuárias. Muitos produtores ainda utilizam métodos tradicionais e manuais, o que dificulta o controle preciso do histórico dos animais, da produção de leite, do estoque de insumos e do controle sanitário. A falta de dados estruturados e a baixa familiaridade com tecnologias digitais limitam a eficiência da gestão e comprometem a sustentabilidade financeira das propriedades. Além disso, o modelo atual do agronegócio brasileiro apresenta concentração de insumos e crédito em grandes empresas, dificultando o acesso da agricultura familiar a recursos e assistência técnica adequada, o que impacta diretamente na produtividade e no desenvolvimento dessas propriedades.

\section{Tecnologias Presentes no Campo}
\label{sec:fundamentacao-teorica-tecnologias-campo}

A chamada Agro 4.0 representa um conjunto de tecnologias digitais que vêm transformando o campo, promovendo maior precisão, eficiência e sustentabilidade na produção agropecuária. Entre as tecnologias mais utilizadas estão drones, sensores, imagens de satélite, máquinas inteligentes e sistemas de monitoramento que permitem o acompanhamento em tempo real das condições ambientais e produtivas. Essas inovações possibilitam a redução dos custos operacionais e a otimização do uso de insumos, além de melhorar a tomada de decisão dos produtores. Contudo, a adoção dessas tecnologias ainda é mais comum em grandes propriedades, enquanto pequenos e médios produtores enfrentam barreiras como o custo, a complexidade dos sistemas e a falta de capacitação.

\section{Avanço das Tecnologias no Agronegócio e na Pecuária}
\label{sec:fundamentacao-teorica-avanço-tecnologico-agro}

O avanço tecnológico no agronegócio brasileiro tem sido impulsionado pela adoção crescente da agricultura de precisão, \gls{IOT}, automação, robótica e plataformas digitais. A agricultura de precisão, por exemplo, emprega \gls{GPS}, sensores e \textit{big data} para monitorar variáveis do campo com alta exatidão, podendo elevar a produtividade em até 20\% e reduzir custos em até 30\%. A \gls{IOT} conecta dispositivos e sensores que coletam dados em tempo real, facilitando o manejo do solo, clima e saúde animal. A automação e a robótica também têm avançado, com máquinas autônomas que otimizam o plantio e a colheita. Além disso, as plataformas digitais aproximam produtores e mercados, aumentando a eficiência comercial, enquanto tecnologias verdes buscam práticas agroecológicas e conservação ambiental.

Na pecuária, surgiram diversos sistemas de gestão com funcionalidades como registro do histórico dos animais, controle de vacinação, monitoramento da produção e gestão de estoque. Esses softwares organizam as informações e apoiam decisões estratégicas. Entretanto, muitos são voltados para grandes propriedades e apresentam interfaces complexas, dificultando a adoção por pequenos e médios produtores. A carência de soluções acessíveis, intuitivas e adaptadas à agricultura familiar ressalta a necessidade de ferramentas específicas que apoiem a gestão e a sustentabilidade das propriedades.

